\documentclass[10pt,a4paper]{amsart}
\usepackage[margin=19mm]{geometry}
\usepackage{amsmath,amssymb,mathtools}
\usepackage[scr=rsfs,cal=euler]{mathalfa}
\usepackage{xfrac}
\usepackage[unicode,hidelinks,bookmarks=false]{hyperref}
\newcommand{\ie}{i.e.\ }
\newcommand{\cf}{cf.\ }
\newcommand{\resp}{respectively\ }
\newcommand{\Subsection}{\subsection}
\providecommand{\nobreakdash}{\nobreak\hbox{-}}
\setlength{\emergencystretch}{2em}
\multlinegap0pt
\usepackage{bm}
\usepackage{bbm}
\usepackage{dsfont}
\usepackage{xspace}
\usepackage{cancel}
\usepackage{mathtools}
\usepackage{cleveref}
\usepackage{amsthm}
\makeatletter
\@ifundefined{corollary}{\newtheorem{corollary}{Corollary}}{}
\@ifundefined{theorem}{\newtheorem{theorem}{Theorem}}{}
\makeatother
\newcommand{\R}{\mathbb{R}}
\title[Improved regularity for a nonlocal dead-core problem]{Improved regularity for a nonlocal dead-core problem}
\author{Disson dos Prazeres, Rafayel Teymurazyan, Jos\'e Miguel Urbano}
\date{Source: arXiv:2504.09557v2 (CC BY 4.0)}
\begin{document}
\maketitle

\noindent\emph{Source and licence.} Improved regularity for a nonlocal dead-core problem. Version arXiv:2504.09557v2 (CC BY 4.0), distributed under the Creative Commons Attribution 4.0 International licence, as recorded in the arXiv licence metadata for this exact version and shown on the abstract page https://arxiv.org/abs/2504.09557v2. Full rights notice: licenses/arxiv-2504.09557-CC-BY-4.0.txt.\par\smallskip

\noindent\textit{Editorial scope.} Counted unit: the Theorem whose source label is \texttt{gradientestimate} in the article (source0.tex lines 596--602, source label \texttt{gradientestimate}), delivered together with the article's own complete original proof of it (source0.tex lines 604--677, one \texttt{proof} environment of 74 lines). No mathematics is written, repaired or re-derived: statement and proof are the article's own text line for line. Required context retained uncounted: equation (lines 108--110); CNtheorem (lines 224--234); SeminormGrowth (lines 248--261); main result 1 (lines 394--400); c4.1 (lines 519--524). Every definition, display and label of those lines is retained unchanged. Omitted from this package and not redistributed: 1--107, 111--223, 235--247, 262--393, 401--518, 525--595, 603--603, 678--895. Nothing inside a retained excerpt is omitted or altered: every retained body is delivered exactly as the article's lines are written, so no omission receipt of any kind accompanies this package. Citations: the retained excerpts carry no citation command at all, so no bibliography entry of the article is transcribed and no reference is redistributed. Reference adaptation: none was needed and none was made; every reference inside the delivered unit resolves inside it, so no unresolved reference or citation is produced. Printed slips kept literal: no printed word, symbol, hypothesis or number of the counted statement or of its proof was altered. Layout and class: the article's own class is \texttt{amsart} with packages including amsmath, amssymb, amsthm, epsfig, hyperref, ulem, euscript, inputenc, xcolor, cleveref, enumerate, mathrsfs, pstricks, pst-grad; the wrapper is the lane's pinned \texttt{amsart} editorial wrapper at 10pt on A4, which restates the theorem-like environments the retained text needs and adds the article's own macro definitions that the retained text needs (\texttt{\textbackslash R}), with extra wrapper packages (bm, bbm, dsfont, xspace, cancel, mathtools, cleveref). The delivered numbering is the unit's own. Assets: no retained span contains a figure environment, a graphics-inclusion command, a PDF-page inclusion, a table or a TikZ picture; no image file is copied into this package, the package declares no asset dependency and no image file is redistributed with it. Licence: version arXiv:2504.09557v2 of this article is distributed under the Creative Commons Attribution 4.0 International licence, as recorded in the arXiv licence metadata for this exact version and shown on the abstract page https://arxiv.org/abs/2504.09557v2. A full rights notice is delivered at licenses/arxiv-2504.09557-CC-BY-4.0.txt. Characters: every retained excerpt is pure ASCII, so the delivered engine, XeLaTeX with its default Latin Modern text font, reports no missing character.
\begin{equation}\label{equation}
   -(-\Delta)^su=u_+^\gamma-u_-^\gamma \quad \textrm{in} \ B_1,
\end{equation}

\begin{theorem}\label{CNtheorem}
If $s\in\left(1/2,1\right)$, and $u$ is a bounded solution of
$$
-(-\Delta)^su=f \quad \textrm{in} \ B_1,
$$
where $f\in L^\infty(B_1)$, then $u\in C^{1+\sigma}(B_{1/2})$, for any $\sigma<2s-1$. Moreover,
$$
\|u\|_{C^{1+\sigma}(B_{1/2})}\le C\left(\|u\|_{L^\infty(\R^n)}+\|f\|_{L^\infty(B_1)}\right),
$$
for a constant $C>0$, depending only on $n$.
\end{theorem}

\begin{theorem}\label{SeminormGrowth}
    Let $s\in(1/2,1)$ and
    \begin{equation}\label{Zeroquotient}
    (-\Delta)^s \left( u(x+h)-u(x) \right)=0, \quad x\in B_1,  
    \end{equation}    
    for all $h\in\R^n$. If, for $0\le\alpha<\beta<1$, there holds
    \begin{equation}\label{crescimento da norma}
    [u]_{C^{\nu+\alpha}(B_R)}\le CR^{\beta-\alpha},\quad \forall R\ge1,
    \end{equation}
    where $\nu=1,2$, and $C>0$ is a constant, depending only on $\alpha$ and $\beta$, then 
    $$
    u(x)=u(0)+Du(0) \cdot x+ \frac{\nu-1}{2} x^T \cdot D^2u(0) x, \quad x\in\R^n.
    $$
\end{theorem}

\begin{theorem}\label{main result 1}
    Let $\gamma\in(0,1/3)$ and $s>1-\frac{\gamma}{2}$. If $u\in L^\infty(\R^n)$ is a viscosity solution of \eqref{equation} and $x_0\in B_{1/2}$ is its branching point, then there exists a constant $C>0$, depending only on $n$ and $\gamma$, such that
\begin{equation}\label{growth estimate 2025}  
    |u(x)|\le C\|u\|_{L^\infty(\R^n)}|x-x_0|^{\frac{2s}{1-\gamma}}, \quad \forall x \in B_{1/2}(x_0).      
\end{equation} 
As a consequence, $u\in C^{\frac{2s}{1-\gamma}}$ at branching points.
\end{theorem}

\begin{corollary}\label{c4.1}
     Under the conditions of \Cref{main result 1}, there exists a constant $C>0$, depending only on $n$ and $\gamma$, such that for any $r<1/2$ one has
    $$
     \sup_{B_r(x_0)}|u|\leq  C\|u\|_{L^\infty(\R^n)}r^{\frac{2s}{1-\gamma}}.
    $$   
\end{corollary}

\begin{theorem}
Under the conditions of \Cref{main result 1}, there exists $C>0$, depending only on $n$ and $\gamma$, such that
\begin{equation}\label{gradientestimate}
    \sup_{B_r(x_0)}|Du|\le C\|u\|_{L^\infty(\R^n)}r^{\frac{2s}{1-\gamma}-1},
\end{equation}
for $r<1/2$, where $x_0\in B_{1/2}$ is a branching point.
\end{theorem}

\begin{proof}    
    Without loss of generality, let $x_0=0$ and $\|u\|_{L^\infty(\R^n)}=1$. Set 
    $$
    \beta:=\frac{2s}{1-\gamma}-1.
    $$
    We argue by contradiction, assuming \eqref{gradientestimate} fails. Then, for each $k\in\mathbb{N}$, there exists $u_k$, a normalized solution of \eqref{equation}, such that, for some $r<1/2$, one has
    $$
    \sup_{B_{r}}|Du_k|>k\,r^\beta.
    $$    
    Let $k_r\in\mathbb{N}$ be such that
    $$
    2^{-(k_r+1)}\le r<2^{-k_r}.
    $$
    Then
    $$
    \sup_{B_{2^{-k_r}}}|Du_k|\ge\sup_{B_r}|Du_k|>k\,r^\beta\ge k2^{-(k_r+1)\beta},
    $$
    \textit{i.e.},    
    \begin{equation}\label{contradictoryassumption}
        \mu_{k_r}>k2^{-(k_r+1)\beta}, 
    \end{equation}
    where
    $$
    \mu_{k_r}:=\sup_{B_{2^{-k_r}}}|Du_k|.
    $$    
    By \Cref{c4.1}, there exists a constant $C>0$, depending only on $n$ and $\gamma$, such that  
    \begin{equation}\label{supestimate}      
    \sup_{B_1}|u_k(2^{-k_r} x)|=\sup_{B_{2^{-k_r}}}|u_k(x)|\le C2^{-k_r(\beta+1)}.
    \end{equation}   
    Set      
    $$
    v_k(x):=\frac{u_k(2^{-k_r}x)}{2^{-k_r}\mu_{k_r}}, \quad x\in B_{1}.
    $$
    Employing \eqref{supestimate} and \eqref{contradictoryassumption}, we estimate
    \begin{equation}\label{vk bound}
    |v_k(x)|\le C\frac{2^{-k_r(\beta+1)}}{2^{-k_r}k2^{-(k_r+1)\beta}}=C\frac{2^\beta}{k}.
    \end{equation}    
    Note that
    \begin{equation}\label{sup of grad vk}      \sup_{B_{1}}|Dv_k|=\sup_{B_{2^{-k_r}}}\frac{|Du_k(x)|}{\mu_{k_r}}=1.
    \end{equation}    
   Furthermore, from \eqref{contradictoryassumption} and \eqref{vk bound}, we get
   \begin{align}\label{fractional Laplacian estimate}
    |(-\Delta)^{s}v_k(x)|&=\mu_{k_r}^{-1}2^{(1-2s)k_r}|(-\Delta)^s u_k(2^{-k_r}x)|\nonumber\\
    &=\mu_{k_r}^{-1}2^{(1-2s)k_r}\left|(u_k)_+^{\gamma}(2^{-k_r}x)-(u_k)_-^\gamma(2^{-k_r}x)\right|\nonumber\\
    &=\mu_{k_r}^{\gamma-1}2^{(1-2s-\gamma)k_r}\left|(v_k)_+^{\gamma}(x)-(v_k)_-^\gamma(x)\right|\nonumber\\
    &<\left[k2^{-(k_r+1)\beta}\right]^{\gamma-1}2^{(1-2s-\gamma)k_r}\left|(v_k)_+^{\gamma}(x)-(v_k)_-^\gamma(x)\right|\nonumber\\
    &=k^{\gamma-1}2^{2s+\gamma-1}\left|(v_k)_+^{\gamma}(x)-(v_k)_-^\gamma(x)\right|\nonumber\\
    &\le\frac{C^{\gamma}2^{2s+\beta\gamma+\gamma}}{k}.
    \end{align}    
    Additionally, for any $0<\alpha<\sigma<2s-1$ and $2^{k_r-1}\ge R\ge1$, recalling \eqref{contradictoryassumption} and \Cref{CNtheorem}, one has
    \begin{align}\label{vk seminorm bound}
        [v_k]_{C^{1+\alpha}(B_R)}&=\frac{2^{-k_r(1+\alpha)}}{\mu_{k_r}}[u_k]_{C^{1+\alpha}(B_{2^{-k_r}R})}\nonumber\\
        &\le\frac{2^{-k_r(1+\alpha)}}{\mu_{k_r}}(2^{-k_r+1}R)^{\sigma-\alpha}[u_k]_{C^{1+\sigma}(B_{2^{-k_r}R})}\nonumber\\   
        &<\frac{2^{-k_r(1+\alpha)+(-k_r+1)(\sigma-\alpha)}}{k2^{-(k_r+1)\beta}}R^{\sigma-\alpha}[u_k]_{C^{1+\sigma}(B_{2^{-k}R})}\nonumber\\
        &=\frac{2^{\sigma-\alpha+\beta}}{k}2^{k_r\left(\frac{2s}{1-\gamma}-2-\sigma\right)}R^{\sigma-\alpha}[u_k]_{C^{1+\sigma}(B_{2^{-k}R})}\nonumber\\
        &\le\tilde{C}R^{\sigma-\alpha},
    \end{align}
    for a constant $\tilde{C}>0$, depending only on $\gamma$ and $s$. Here, the last inequality follows by choosing $\sigma>0$ sufficiently close to $2s-1$ so
    $$
    \frac{2s}{1-\gamma}-2-\sigma<0.
    $$
    Thus, up to a subsequence, $v_k$ converges to some $v_0$ in $C^{1+\alpha}(B_R)$, as $k\to\infty$. Moreover, thanks to \eqref{fractional Laplacian estimate} and \eqref{vk seminorm bound},
    $$
    (-\Delta)^sv_0=0
    $$
    and
    $$
    [v_0]_{C^{1+\alpha}(B_R)}\le\tilde{C}R^{\sigma-\alpha}.
    $$
    \Cref{SeminormGrowth} applies to $v_0$, yielding $v_0\equiv0$, which is a contradiction, since from \eqref{sup of grad vk}, one has
    $$
    \sup_{B_{1}}|Dv_0|=1.
    $$
\end{proof}

\end{document}
